\documentclass[11pt,a4paper]{article}

\usepackage[a4paper,margin=1in]{geometry}
\usepackage[T1]{fontenc}
\usepackage[utf8]{inputenc}
\usepackage{lmodern}
\usepackage{xcolor}
\usepackage{graphicx}
\usepackage{hyperref}
\usepackage{listings}
\usepackage{enumitem}
\usepackage{array}
\usepackage[most]{tcolorbox}

\definecolor{DeloitteGreen}{HTML}{86BC25}
\definecolor{DeloitteGreenDark}{HTML}{1F3B08}
\definecolor{DeloitteInk}{HTML}{171717}
\definecolor{DeloitteSoft}{HTML}{F5F8EE}
\definecolor{DeloitteBlueSoft}{HTML}{EEF6FB}
\definecolor{DeloitteBlueBorder}{HTML}{6FA8DC}

\hypersetup{
    colorlinks=true,
    linkcolor=DeloitteGreenDark,
    urlcolor=DeloitteGreenDark,
    pdftitle={Deloitte SQL Agent Installation Guide},
    pdfauthor={OpenAI Codex}
}

\graphicspath{{./img/}{img/}{./docs/img/}{docs/img/}}

\lstset{
    basicstyle=\ttfamily\small,
    breaklines=true,
    columns=fullflexible,
    frame=single,
    framerule=0.4pt,
    rulecolor=\color{black!30},
    backgroundcolor=\color{black!2}
}

\newtcolorbox{infobox}{
    enhanced,
    breakable,
    colback=DeloitteBlueSoft,
    colframe=DeloitteBlueBorder,
    boxrule=0.6pt,
    arc=3pt,
    left=10pt,
    right=10pt,
    top=8pt,
    bottom=8pt
}

\newtcolorbox{titlepanel}{
    enhanced,
    colback=DeloitteInk,
    colframe=DeloitteInk,
    boxrule=0pt,
    arc=0pt,
    left=18pt,
    right=18pt,
    top=20pt,
    bottom=20pt
}

\newcommand{\screenshotplaceholder}[2][4.8cm]{%
    \begin{figure}[ht]
    \centering
    \fbox{%
        \parbox[c][#1][c]{0.9\linewidth}{%
            \centering
            \textbf{Screenshot Placeholder}\\[0.4em]
            #2
        }%
    }
    \caption{#2}
    \end{figure}
}

\newcommand{\screenshot}[3][0.92\linewidth]{%
    \begin{figure}[ht]
    \centering
    \includegraphics[width=#1]{#2}
    \caption{#3}
    \end{figure}
}

\newcommand{\guidecover}{%
    \begin{titlepanel}
    {\color{DeloitteGreen}\rule{\linewidth}{1.5pt}}\\[1.2em]
    {\Huge\bfseries\color{white}Deloitte SQL Agent}\\[0.35em]
    {\LARGE\color{white}Installation Guide}\\[1em]
    {\large\color{white!85}Windows 11 VM setup via SharePoint ZIP distribution}\\[2.2em]
    {\color{white!88}
    \begin{tabular}{>{\bfseries}p{3.2cm}p{9.5cm}}
    Audience & Deloitte users installing the agent on the approved Windows 11 VM \\
    Platform & Windows 11, PowerShell, OpenRouter, Python 3.14, uv \\
    Version date & April 2, 2026 \\
    \end{tabular}
    }
    \end{titlepanel}
    \vspace{1.2em}
}

\begin{document}
\guidecover

\section{Purpose}

This document explains how to install and run the Deloitte SQL Agent on a Deloitte-provided Windows 11 virtual machine. The agent is distributed as a ZIP file via SharePoint.

\section{Prerequisites}

Before you begin:

\begin{itemize}[leftmargin=1.5em]
    \item Use the Deloitte Windows 11 VM.
    \item Connect the VM to \textbf{Telekom WLAN}, \textbf{not DeloitteNet}.
    \item Make sure you can download files from SharePoint and browse to public websites.
    \item This guide assumes PowerShell is available.
\end{itemize}

This guide does not explain how to provision the VM itself. Use the Deloitte VM onboarding tutorial:

\begin{quote}
\href{https://deloitteemea.service-now.com/mysupport?id=sc_cat_item&sys_id=4e02d262c3a7c2145608282ce00131bf&table=sc_cat_item&searchTerm=VMware%20Workstation%20Player}{Deloitte VM onboarding tutorial}
\end{quote}

\section{Step 1: Install uv}

uv is the Python package manager used by this project. On Deloitte machines, prefer the official WinGet package.\footnote{\url{https://docs.astral.sh/uv/getting-started/installation/}}

Open \textbf{PowerShell} and run:

\begin{lstlisting}
winget install --id=astral-sh.uv -e
\end{lstlisting}

Close PowerShell and open it again, then verify the installation:

\begin{lstlisting}
uv --version
\end{lstlisting}

If WinGet is unavailable, you can use the official Astral PowerShell installer as an alternative. At the time this guide was written, this installer was often blocked in the Deloitte setup even on Telekom WLAN, so it should not be the primary path:

\begin{lstlisting}
powershell -ExecutionPolicy ByPass -c "irm https://astral.sh/uv/install.ps1 | iex"
\end{lstlisting}

\screenshot{install_uv.png}{uv installation via WinGet and the verification step in PowerShell.}

\section{Step 2: Install Python 3.14 with uv}

This project requires Python 3.14 or newer. The requirement is defined in the project metadata.

Install Python 3.14 from PowerShell:

\begin{lstlisting}
uv --native-tls python install 3.14
\end{lstlisting}

\begin{infobox}
\textbf{Why \texttt{--native-tls}?}

Use \texttt{--native-tls} on Deloitte VMs so \texttt{uv} reads certificates from the Windows certificate store. This is the recommended way to make \texttt{uv} work in managed environments with custom trust settings.\footnote{\url{https://docs.astral.sh/uv/reference/cli/}}
\end{infobox}

You can verify that uv sees the installed interpreter:

\begin{lstlisting}
uv --native-tls python list
\end{lstlisting}

Reference: official uv Python installation guide.\footnote{\url{https://docs.astral.sh/uv/guides/install-python/}}

\screenshot{install_python.png}{Python 3.14 installation with \texttt{uv --native-tls python install 3.14}.}

\section{Step 3: Download and Extract the ZIP from SharePoint}

\begin{enumerate}[leftmargin=1.5em]
    \item Open the SharePoint link provided by the Deloitte team.
    \item Download the ZIP file, for example \texttt{deloitte-sql-agent.zip}.
    \item Right-click the ZIP file and choose \textbf{Extract All}.
    \item Extract it to a simple working folder, for example:
\end{enumerate}

\begin{quote}
\texttt{C:\textbackslash Users\textbackslash <your-user>\textbackslash Projects\textbackslash deloitte-sql-agent}
\end{quote}

\screenshotplaceholder{Add a screenshot of the SharePoint download page and extracted folder.}

\section{Step 4: Install the Project Dependencies}

Open PowerShell in the extracted project folder and run:

\begin{lstlisting}
cd C:\Users\<your-user>\Projects\deloitte-sql-agent
uv --native-tls sync
\end{lstlisting}

This creates the local virtual environment and installs the project dependencies from the lockfile.

\screenshot[0.72\linewidth]{uv_sync.png}{Successful dependency installation with \texttt{uv --native-tls sync}.}

\section{Step 5: Create an OpenRouter Token}

The SQL Agent calls OpenRouter for model inference, so each user needs an OpenRouter API key (token).

\subsection*{5.1 Sign in}

Open the OpenRouter sign-in page:

\begin{quote}
\url{https://openrouter.ai/sign-in}
\end{quote}

\subsection*{5.2 Set up payment if required}

Depending on your OpenRouter account or workspace, you may need to add payment information before you can create and use an API key.

\begin{enumerate}[leftmargin=1.5em]
    \item Open your OpenRouter account or workspace settings.
    \item Add payment details if required by your setup.
    \item Return to the API key area after payment is configured.
\end{enumerate}

\subsection*{5.3 Open the keys page}

After signing in, open the keys page:

\begin{quote}
\url{https://openrouter.ai/settings/keys}
\end{quote}

At the time this document was written on April 2, 2026, this URL redirects to the workspace key-management page after login.

\subsection*{5.4 Create and copy the token}

\begin{enumerate}[leftmargin=1.5em]
    \item Click the button to create a new API key.
    \item Give the key a clear name, for example \texttt{Deloitte SQL Agent}.
    \item Copy the new key immediately.
\end{enumerate}

\begin{infobox}
\textbf{OpenRouter user interface note}

The OpenRouter website may change over time. If buttons or menu labels look slightly different, the overall process should still be similar: sign in, set up payment if needed, then create an API key.
\end{infobox}

OpenRouter references:

\begin{itemize}[leftmargin=1.5em]
    \item \href{https://openrouter.ai/docs/quickstart}{OpenRouter Quickstart}
    \item \href{https://openrouter.ai/docs/api-reference/api-keys/get-current-api-key}{OpenRouter API key documentation}
\end{itemize}

\textbf{Important:} Treat the token like a password. Do not send it by email or save it in screenshots.

\screenshot{open_router.png}{OpenRouter during payment and token setup.}

\section{Step 6: Run the SQL Agent}

Run the agent from the project folder. The example below uses a SQLite database file.

\begin{lstlisting}
uv --native-tls run python main.py `
  --db-url sqlite:///C:/path/to/source.db `
  --target-rows 1000 `
  --out subset.sql `
  --model moonshotai/kimi-k2-thinking `
  --api-key sk-or-v1-REPLACE_ME `
  --no-verify-ssl
\end{lstlisting}

\begin{infobox}
\textbf{Why \texttt{--no-verify-ssl}?}

The SQL Agent itself uses Python HTTP requests for the OpenRouter call. In the Deloitte VM setup, run the agent with \texttt{--no-verify-ssl}. Only use this on the approved VM and trusted Telekom WLAN setup.
\end{infobox}

Notes:

\begin{itemize}[leftmargin=1.5em]
    \item \texttt{uv --native-tls} is used for the project launcher.
    \item \texttt{--no-verify-ssl} is passed to the SQL Agent itself.
    \item The tutorial uses \texttt{--api-key} directly on the command line.
    \item The output file \texttt{subset.sql} will be written into the current directory unless you provide a different path.
\end{itemize}

\begin{infobox}
\textbf{Optional alternative: environment variable}

If you do not want to pass the key on the command line every time, you can set \texttt{OPENROUTER\_API\_KEY} in PowerShell and omit \texttt{--api-key} from the run command.
\end{infobox}

\screenshot{agent_running.png}{Successful first run of the SQL Agent in PowerShell.}

\section{Troubleshooting}

\subsection*{uv command is not found}

Close PowerShell and open it again. If that still does not help, reinstall uv and verify with:

\begin{lstlisting}
uv --version
\end{lstlisting}

\subsection*{TLS or certificate errors during uv commands}

\begin{itemize}[leftmargin=1.5em]
    \item Confirm that the VM is connected to \textbf{Telekom WLAN}, not DeloitteNet.
    \item Make sure the command includes \texttt{--native-tls}.
    \item If the issue continues, contact the Deloitte support contact for the VM environment.
\end{itemize}

\subsection*{TLS or certificate errors during the SQL Agent run}

\begin{itemize}[leftmargin=1.5em]
    \item Confirm that the run command includes \texttt{--no-verify-ssl}.
    \item Confirm that the VM is using \textbf{Telekom WLAN}.
    \item Retry the command in a fresh PowerShell window.
\end{itemize}

\subsection*{OpenRouter authentication failed}

\begin{itemize}[leftmargin=1.5em]
    \item Check that the value passed to \texttt{--api-key} is correct.
    \item If you use \texttt{OPENROUTER\_API\_KEY} instead, verify that the variable is set correctly.
    \item Generate a new token in OpenRouter if necessary.
\end{itemize}

\section{Quick Reference}

\begin{lstlisting}
winget install --id=astral-sh.uv -e
uv --native-tls python install 3.14
cd C:\Users\<your-user>\Projects\deloitte-sql-agent
uv --native-tls sync
uv --native-tls run python main.py `
  --db-url sqlite:///C:/path/to/source.db `
  --target-rows 1000 `
  --out subset.sql `
  --model moonshotai/kimi-k2-thinking `
  --api-key sk-or-v1-REPLACE_ME `
  --no-verify-ssl
\end{lstlisting}

\end{document}
